\head{Dot 续译整合核查：预备笔记第12--13节}
\source{\texttt{EGAVaiello.pdf}，PDF及印刷第43--53页；Dot 续作工程包中的19批实际TeX源码}
本节说明这19批译文在本次合订时的处理。连续范围从第12节标题及记号补充起，至第13节注（13.4）e）末止，下一原文标题为第14节“锥投影”。各批的封面、独立页码和旧接续说明已退出正文；原文段落、公式、英译编者注以及译校脚注均纳入相应位置。完整原批文件及原有\texttt{AUDIT.txt}仍保存在工程档案中。

这次核对了19份实际TeX、对应的进度记录和底本第43--53页文本，并查看了第43、44、47、50、51、53页图像，抽核秩及对偶约定、退化层维数公式、引理（13.3）的假设、一般纤维基域，以及（13.4）的三项附注。旧批声称逐页检查、编译成功的记录属于历史记录，不作为本次重新执行同样检查的证据。本次全书编译及页面检查另见总技术核查。

\subsec{按原文顺序的入编清单}
下表简称与工程中的完整文件名、SHA-256及旧进度字段逐项对应，详见\texttt{integration\_notes.json}。页码均指本次英文编辑底本，恰与该文件显示的印刷页码相同。
\begin{longtable}{p{.08\textwidth}p{.57\textwidth}p{.20\textwidth}}
\toprule 序号 & 译文范围 & 底本页 \\\midrule\endhead
1 & 第12节开篇：记号及关联预概形 & 43--44\\
2 & 记号自评、格拉斯曼纤维及射影丛 & 44\\
3 & 对第2--4节的一般截面推广指示 & 44--45\\
4 & 对第5节正则截面的推广指示 & 45\\
5 & 第6--8节维数条件及推广计划 & 46\\
6 & 奇异关联簇与交维数过滤 & 46--47\\
7 & 退化层维数与超奇异点 & 47\\
8 & 线性截面族、爆破与判别集合 & 48\\
9 & 第12节末待证提纲；定义（13.1） & 48--49\\
10 & 定理（13.2）及初步归约 & 49\\
11 & 闭边界选择；引理（13.3）陈述 & 50\\
12 & 由引理构造初等邻域 & 50--51\\
13 & 引理（13.3）证明的归约 & 51\\
14 & 超平面补记及对余维的归纳 & 51--52\\
15 & 超平面情形的证明收尾 & 52\\
16 & 证明后附记；注（13.4）a）、b） & 52--53\\
17 & 注（13.4）c）：有限点集邻域 & 53\\
18 & 注（13.4）d）：一般基底 & 53\\
19 & 注（13.4）e）：拓扑结构 & 53\\
\bottomrule
\end{longtable}

\subsec{结构保全与本次改动}
19批共有93条\texttt{note}译注，已逐条保留并改成\texttt{D17-M001}至\texttt{D17-M093}的脚注编号；六条底本脚注61--66保持，并明确其英译编者来源。第（13.1）--（13.4）四个编号条目现使用与合订本一致的条目宏，保留原编号。各批为避免孤页设置的强制换页、加长页面指令已删除，由\texttt{ctexbook}统一分页。

术语作三项统一：准凝聚改为拟凝聚，未分歧改为无分歧，多初等改为多重初等。旧注中“本批止于这里”“后续构造未完成”等因分批而产生的状态文字已改为实际合订状态；原文自身的待证、难辨及推测语气保留。本次另行订正一般纤维的基域，以及退化层维数公式的非空适用界；后者有具体反例和线性代数维数计算为据。这些局部核验并不构成对整套数学论证的重新验证。

\audit{D17-M-E01}{一般线性截面的光滑基域}{英文原句的基域措辞不适用于其所指一般纤维；正文订正，依赖链仍待核}
底本第51页在令$\eta$为参数子格拉斯曼簇的一般点之后，写相应截面“smooth over $k$”；第52页随后明确使用$X^{(m)}_\eta$。该一般纤维的结构基域是$\kappa(\eta)$。本次正文改为“$X^{(m)}_\eta$在$\kappa(\eta)$上光滑”，并在脚注\texttt{D17-M065}保留原写法和理由。

这不是可以任意忽略的基域名称差异。光滑态射须局部有限表示；可参照\href{https://stacks.math.columbia.edu/tag/01V4}{The Stacks Project，Tag 01V4，光滑态射定义及局部有限表示引理}。例如，$\Spec k(t)\to\Spec k$并非局部有限表示：若$k(t)$可由有限个分式作为$k$-代数生成，取其分母乘积$d(t)$，生成的代数将包含在$k[t,1/d(t)]$中。选一个不整除$d$的不可约多项式$q(t)$，则$1/q(t)$不在这个环中，与它应等于$k(t)$矛盾。由此不能仅因一个一般纤维相对于$\kappa(\eta)$光滑，就称它对原域$k$光滑。

这里仅订正一般纤维应使用的基域，未由这项订正证明草稿中的Bertini断言。$Q_n$、$\Phi_n$、$\varphi_n$三种符号及8.8缺文的依赖仍照具体位置分别提示。

\audit{D17-M-E02}{退化层维数公式的参数限制}{已订正非空适用范围，并明确几何纤维纯维约定；原错误界及计算语气存于脚注}
第47页对$V^{(k)}$提出
\[
\dim\operatorname{Grass}^{m}(P)-\dim V^{(k)}=(k-1)(n-m)+k^2,
\]
给出的限制是$k\leq m$、$k\leq r-m$，并写“除非计算有误”。这两项限制不能保证退化层非空。例如取$S=\Spec K$、$X=P=\mathbf P^2_K$、$f=\operatorname{id}$，并取$r=3$、$n=2$、$m=k=1$。此时$X$光滑、$f$无分歧，分支维数$n$满足$n\geq m$，两项原限制也均成立。对每个关联点$(x,L)$，$T_x=P$而$L$是直线，故
\[
\dim(T_x\cap L)=1<2=n-m+k.
\]
所以$V^{(1)}=\varnothing$，原式所写的有限维数计算不能在这里照用。

本次将计算明确限定在一个非空、纯维$n$的几何纤维上；此时$P$的维数为$r-1$，原文已有$n\geq m$，无分歧性又给出$n\leq r-1$。在该纤维内，正确的非空参数范围为
\[
0\leq m\leq n\leq r-1,\qquad
0\leq k\leq\min(m,r-1-n).
\]
正文据此订正，脚注\texttt{D17-M034}保存原限制与反例。下面给出本次译校新增的线性代数验证；这段验证不是原稿省略证明的假托恢复。

固定一个几何点$x$，令$N=r-1$、$d=N-m$。经过$f(x)$的余维$m$线性子簇$L$，对应于$N$维切向量空间$V$中的$d$维子空间$W$。无分歧性使$T_xX$作为一个$n$维子空间$A\subset V$出现。原交维条件等价于
\[
\dim(A\cap W)\geq t,\qquad t=n-m+k.
\]
交维恰为$t$的层可以依次构造：选择$J=A\cap W$（$A$中$t$维子空间），选择$B$（$W$在$V/A$中的$(d-t)$维像），再选择$B\to A/J$的图映射。因此其维数是
\begin{align*}
&t(n-t)+(d-t)\bigl((N-n)-(d-t)\bigr)+(d-t)(n-t)\\
&\hspace{2em}=t(n-t)+(d-t)m
=dm-k(n-m+k).
\end{align*}
这些选择非空，当且仅当$\max(0,n-m)\leq t\leq\min(n,d)$；在$n\geq m$的假设下，正好化为上列$k$的范围。将$k$增加一，余维$k(n-m+k)$增加$n-m+2k+1>0$；故交维更大的层维数更小，闭条件$\dim(A\cap W)\geq t$的维数就是所算的数值。

当$x$变化时，相关切空间包含构成向量丛的子丛：$f^*\Omega_{P/K}\twoheadrightarrow\Omega_{X/K}$是到局部自由层的满射，局部可分裂，所以其对偶是局部可分裂的包含。局部选取适配子丛的标架，以上交维族成为相同Schubert层的局部乘积。因此
\[
\dim V^{(k)}=n+m(N-m)-k(n-m+k).
\]
另一方面，参数化$P$中全部余维$m$线性子簇的Grassmann簇维数为$m(N+1-m)$。两式相减得到
\[
m-n+k(n-m+k)=(k-1)(n-m)+k^2,
\]
即原稿中的维数公式。订正的是其适用界和未明说的纤维、纯维约定，而非公式右端。

这项计算给出退化层的维数，不识别其非光滑子概形，更没有补齐补篇第16段的引理或后续超奇异点结论。若几何纤维有不同维数的分支，须在各纯维分支上分别使用相应$n$；不能仍把同一个$n$用于整个混合维数纤维。

\audit{D17-M-C01}{关联预概形的秩与对偶约定}{Dot已写入正文的局部订正得以保留，秩计算相符}
第12节重新约定$E$的秩为$r$，并用
\[
\operatorname{Gr}^{m}(P)=\operatorname{Grass}_m(E^\vee)
\]
参数化相对余维$m$的线性子簇。因此泛商为$E^\vee_{\operatorname{Gr}^{m}}\twoheadrightarrow G$，其对偶才给出$G^\vee\hookrightarrow E$。转置后的箭头目标应为$G_T$，与截面所在的$G_T(1)$一致。这些是由原段自身定义可核对的对象匹配，底本的混用已保留在脚注中。

若$F_X=\ker(E_X\twoheadrightarrow\mathcal O_X(1))$，则$F_X$的秩为$r-1$；所以$\operatorname{Grass}_m(F_X^\vee)$的纤维维数为$m(r-1-m)$，在$0\leq m\leq r-1$时纤维非空且几何整。旧批已把原式$m(r-m)$订正并列出$r=2,m=1$的核验例。这一订正保留；没有把$H^{(m)}\to\operatorname{Gr}^{m}$的射影丛结构混同于$X^{(m)}\to\operatorname{Gr}^{m}$的截面族。

\audit{D17-M-P01}{引理（13.3）的经过指定点条件}{原稿陈述与证明间存在未补齐的依赖，保持待核}
第50页引理（13.3）的陈述没有明列$x\in L_0$，第51页证明却从“经过$x$的一般线性子簇”开始；中间构造欲由所选$L_0$推出所得开集包含$x$，同样使用了这一条件。Dot原注已标出这一遗漏，本次继续保留。

此外，证明归约中写$x\in P(k)$，而定理（13.2）只写$x$为光滑点；注（13.4）c）先将有限点集归约到闭点，又在后句使用$\Phi(k)$。闭点与$k$-有理点并不等价，因此不能仅在引理陈述中增补$x\in L_0$就宣称整条证明已经修复。本次不偷偷加上定理没有列出的有理性假设，也不凭现代理论代替原稿缺少的过渡。相关文字及疑点保持可见。

\audit{D17-M-P02}{第12节推广计划与原稿缺文}{译文连续，相关陈述仍未获得完整证明核查}
第12节大量段落本身就是编写计划和推广指示。4.2/4.3的$m-1$或$m+1$读法、奇异关联簇描述中的残词、过滤层的具体概形结构、补篇引理位置、一般线性族中未写完的$\Spec$对象及第12节最后的几何不可约性引理残句均未据猜测补写。

对于退化层维数公式，旧注只检查了$k=2$代入后为$n-m+4$；本次在D17-M-E02另给具体空轨迹例子，并在非空纯维几何纤维的明确范围内完成线性代数计算；后续关于过滤的概形结构及超奇异点的论证仍未因此得到验证。证明（13.2）/（13.3）援引8.7(c)、8.8，也不能由后文的“因此证明”反推被引条目的缺文已恢复。注（13.4）d）一般基底的条件和e）普遍覆盖同胚于$\mathbb C^n$的较强断言均保留原作者的陈述，同时保留未独立证明的说明。
